\chapter[2007 --- Capecchi et al.]{2007: Mario R. Capecchi, Sir Martin J. Evans, Oliver Smithies}
\label{ch:2007_Capecchi}

\section*{Main Contribution}

\textit{Developed gene targeting in mice using embryonic stem cells, enabling the creation of genetically modified animals to study gene function and model human diseases.}

\section{Official Citation}

for their discoveries of principles for introducing specific gene modifications in mice by the use of embryonic stem cells

\section{Background and Historical Context}

The ability to create genetically modified organisms---animals carrying specific, targeted changes in their DNA---has been one of the most powerful tools in modern biology and medicine. The development of gene targeting technology in mice, for which Mario R. Capecchi, Sir Martin J. Evans, and Oliver Smithies were awarded the 2007 Nobel Prize in Physiology or Medicine, revolutionized the study of gene function and disease mechanisms. Their work enabled researchers to create ``knockout mice''---animals in which a specific gene had been inactivated---providing an notable system for studying the role of individual genes in health and disease.

To appreciate the significance of this achievement, it is necessary to understand the state of genetics and molecular biology in the 1970s and 1980s. While the basic mechanisms of DNA replication, transcription, and translation were understood, the tools available for studying gene function in living organisms were limited. Researchers could introduce foreign DNA into cells, but the integration of that DNA into the genome was random and uncontrolled. The ability to make precise, predetermined changes in a specific gene---to target modifications to a particular location in the genome---was beyond the reach of existing technology.

The concept of homologous recombination---the exchange of DNA sequences between similar or identical DNA molecules---provided a potential solution to this problem. Homologous recombination occurs naturally during meiosis, when homologous chromosomes exchange segments of DNA through a process known as crossing over. If this mechanism could be harnessed in cultured cells, it might be possible to introduce a modified version of a gene that would replace the endogenous copy through homologous recombination, thereby creating a precise genetic modification at a specific genomic location.

The challenge was that homologous recombination between introduced DNA and the chromosomal genome was an extremely rare event. In any given cell, the vast majority of introduced DNA molecules would either fail to integrate or would integrate randomly at non-homologous sites. The frequency of homologous recombination events was estimated to be at least two to three orders of magnitude lower than non-homologous integration, making it extremely difficult to identify the rare cells in which homologous recombination had occurred.

\section{The Discovery/Contribution}

Mario R. Capecchi, born in Italy and educated at the University of Rome and the California Institute of Technology, made foundational contributions to the development of gene targeting technology. In the early 1980s, while at the University of Utah, Capecchi developed methods for introducing DNA into cultured mammalian cells and demonstrated that homologous recombination could occur between introduced DNA and chromosomal sequences. His early work established the basic framework for gene targeting experiments and provided proof of principle that homologous recombination could be used to make specific genetic modifications in mammalian cells.

Capecchi's approach involved constructing a ``targeting vector''---a piece of DNA containing a modified version of the gene of interest, flanked on both sides by sequences homologous to the chromosomal gene. When this vector was introduced into cultured cells, rare homologous recombination events could occur in which the modified gene replaced the endogenous copy. To identify these rare events, Capecchi incorporated a positive selection marker (typically a drug resistance gene) into the targeting vector, along with a negative selection marker positioned outside the region of homology. Cells that underwent homologous recombination would retain the positive marker but lose the negative marker, while cells with random integration would retain both markers. This ``positive-negative selection'' strategy greatly enriched for cells with homologous recombination events and made gene targeting practically feasible.

Oliver Smithies, a British-born scientist working at the University of North Carolina at Chapel Hill, independently developed a similar approach to gene targeting in mammalian cells. Smithies' work on the genetics of the major histocompatibility complex (MHC) and his studies of globin gene regulation provided the biological motivation for developing gene targeting technology. He demonstrated that homologous recombination could be used to introduce specific mutations into the hypoxanthine-guanine phosphoribosyltransferase (HPRT) gene in human cells and showed that these modifications were stably inherited. Smithies also made important contributions to the refinement of targeting vector design and the development of efficient selection strategies.

Sir Martin J. Evans, a British scientist at the University of Cambridge (later Cardiff University), made the critical contribution that linked gene targeting technology in cultured cells to the creation of genetically modified animals. Evans' key innovation was the development of methods for culturing and genetically manipulating embryonic stem (ES) cells---pluripotent cells derived from the inner cell mass of mouse blastocysts. ES cells have the remarkable property of being able to contribute to all tissues of the developing embryo, including the germline, when injected into a host blastocyst.

Evans demonstrated that ES cells could be maintained in culture for extended periods while retaining their pluripotent properties, and that genetically modified ES cells could be used to generate chimeric mice in which the modified cells contributed to various tissues. Crucially, Evans showed that if the modified ES cells contributed to the germline of the chimeric mice, the genetic modification could be passed on to offspring, thereby establishing a line of genetically modified animals. This achievement, first reported in 1981, provided the essential link between gene targeting in cultured cells and the creation of genetically modified animals.

The synthesis of these three contributions---Capecchi and Smithies' gene targeting technology and Evans' ES cell culture and chimera methods---made it possible to create knockout mice with specific genetic modifications. The process involved several steps: first, constructing a targeting vector with the desired genetic modification; second, introducing the vector into cultured ES cells and selecting for homologous recombination events; third, injecting the genetically modified ES cells into host blastocysts to create chimeric mice; and fourth, breeding the chimeric mice to establish lines of animals carrying the genetic modification in all cells, including the germline.

The first knockout mice were produced in the mid-1980s, and the technology was rapidly adopted by researchers worldwide. By the early 1990s, hundreds of genes had been knocked out in mice, providing invaluable information about gene function in development, physiology, and disease. The creation of knockout mice became one of the most widely used experimental approaches in biomedical research, enabling researchers to test hypotheses about gene function with notable precision and rigor.

\section{Impact on Modern Medicine}

The development of gene targeting technology in mice has had a significant impact on virtually every area of biomedical research and has fundamentally changed our understanding of gene function and disease mechanisms. Knockout mice have been used to elucidate the roles of thousands of genes in normal development and physiology, and have provided critical insights into the molecular basis of human diseases.

One of a major applications of knockout mice has been in the study of human genetic diseases. By knocking out the mouse homologs of genes implicated in human diseases, researchers have been able to create animal models that recapitulate key features of the human condition. These models have been invaluable for studying disease mechanisms, testing therapeutic interventions, and developing new treatments. For example, knockout mice for the tumor suppressor genes \textit{p53} and \textit{Rb} have provided fundamental insights into cancer biology, while knockout models for genes involved in cystic fibrosis, sickle cell disease, and Huntington's disease have advanced our understanding of these conditions and facilitated the development of potential therapies.

The technology has also been instrumental in understanding the immune system. Knockout mice have been used to dissect the roles of individual immune system genes and to understand how defects in these genes lead to immunodeficiency, autoimmunity, and allergic diseases. The development of mouse models for HIV/AIDS, using knockout mice lacking the \textit{CCR5} co-receptor, has provided important insights into viral pathogenesis and potential therapeutic strategies.

In neuroscience, knockout mice have been used to study the roles of neurotransmitter receptors, ion channels, and signaling molecules in brain function and behavior. These studies have contributed to our understanding of neurological and psychiatric disorders, including epilepsy, depression, schizophrenia, and neurodegenerative diseases. The ability to create mice with specific genetic modifications in defined brain regions has opened new avenues for studying brain function and developing targeted therapies.

The impact of gene targeting technology extends beyond the creation of simple knockout mice. Refinements of the original technology have enabled the creation of more sophisticated genetic modifications, including conditional knockouts (in which gene inactivation can be induced at specific times or in specific tissues), knock-in mice (in which a specific mutation is introduced without inactivating the gene), and reporter mice (in which the expression of a gene can be visualized in living tissues). These advanced techniques have further expanded the utility of genetically modified mice as experimental tools.

The development of gene targeting technology has also contributed to the advancement of gene therapy approaches for human genetic diseases. The understanding of homologous recombination and gene targeting in mouse cells has informed the development of strategies for correcting disease-causing mutations in human cells, both \textit{ex vivo} and \textit{in vivo}. The recent success of CRISPR-based gene editing therapies, while using different molecular tools, builds upon the conceptual framework established by Capecchi, Evans, and Smithies.

In 2007, Mario R. Capecchi, Sir Martin J. Evans, and Oliver Smithies were awarded the Nobel Prize in Physiology or Medicine ``for their discoveries of principles for introducing specific gene modifications in mice by the use of embryonic stem cells.'' The Nobel Committee recognized the transformative nature of their work and its impact on our understanding of gene function and disease.

\section{Legacy}

The legacy of Capecchi, Evans, and Smithies' work is woven into the fabric of modern biomedical research. Gene targeting technology in mice has become an indispensable tool in virtually every area of biology and medicine, and knockout mice continue to be generated and studied in thousands of laboratories worldwide.

The creation of comprehensive knockout mouse libraries has been a major initiative in genomics research. Projects such as the International Knockout Mouse Consortium (IKMC) and the Knockout Mouse Phenotyping Program (KOMP) have aimed to create and characterize knockout mice for every gene in the mouse genome, providing a comprehensive resource for studying gene function on a genome-wide scale. These efforts have generated thousands of knockout mouse lines, each of which has been phenotyped for a wide range of physiological and behavioral parameters.

The development of gene targeting technology has also inspired the development of new genome editing tools, most notably CRISPR-Cas9, which has revolutionized genetic engineering in virtually all organisms. While CRISPR-Cas9 uses a different molecular mechanism than the homologous recombination-based approach developed by Capecchi, Evans, and Smithies, the conceptual goal of making precise, targeted genetic modifications is the same. The success of gene targeting in mice provided both the biological motivation and the experimental framework for the development of these newer technologies.

The work of Capecchi, Evans, and Smithies also highlighted the importance of basic research and technological innovation in advancing our understanding of disease. Their development of gene targeting technology was driven by fundamental questions about gene function, not by specific clinical applications. Yet the impact of their work on clinical medicine has been immense, illustrating the unpredictable but significant ways in which basic research can lead to therapeutic advances.

Today, the knockout mouse remains one of the major experimental models in biomedical research. New technologies, including CRISPR-based genome editing and single-cell genomics, are being combined with gene targeting to create increasingly sophisticated models of human disease. The legacy of Capecchi, Evans, and Smithies continues to inspire new generations of researchers who use genetically modified mice to explore the mysteries of gene function and to develop new treatments for human diseases.


\section{Modern Developments}

Capecchi, Evans, and Smithies' gene targeting work has led to modern genetic engineering. Knockout and knock-in mice are essential tools for studying gene function and modeling human diseases. Understanding of homologous recombination has informed development of CRISPR-based gene editing. Conditional knockout mice enable tissue-specific and time-specific gene deletion. Humanized mice carrying human genes are used to study human diseases and test therapies.
